\documentclass[11pt,a4paper]{article}
\usepackage[T1]{fontenc}
\usepackage{isabelle,isabellesym}

\usepackage{todonotes}


\usepackage{amsmath,amssymb}

% further packages required for unusual symbols (see also
% isabellesym.sty), use only when needed

%\usepackage{amssymb}
  %for \<leadsto>, \<box>, \<diamond>, \<sqsupset>, \<mho>, \<Join>,
  %\<lhd>, \<lesssim>, \<greatersim>, \<lessapprox>, \<greaterapprox>,
  %\<triangleq>, \<yen>, \<lozenge>

%\usepackage{eurosym}
  %for \<euro>

%\usepackage[only,bigsqcap]{stmaryrd}
  %for \<Sqinter>

%\usepackage{eufrak}
  %for \<AA> ... \<ZZ>, \<aa> ... \<zz> (also included in amssymb)

%\usepackage{textcomp}
  %for \<onequarter>, \<onehalf>, \<threequarters>, \<degree>, \<cent>,
  %\<currency>

% this should be the last package used
\usepackage{pdfsetup}

% urls in roman style, theory text in math-similar italics
\urlstyle{rm}
\isabellestyle{it}

% for uniform font size
%\renewcommand{\isastyle}{\isastyleminor}


\begin{document}

\title{Heterogeneous Dynamic Logic}
%\author{}
\maketitle

\begin{abstract}

\end{abstract}

\section{Preamble}
This Isabelle package formalizes \emph{Dynamic Logic}.
Dynamic Logic can be thought of as a multi-modal logic where we can
parameterize the possible state transitions inside a modality.
The intention of this package is to formalize the fundamental notions that
make a logic a dynamic logic and prove the formulas that hold for a dynamic
logic irrespective of its concrete instantiation.
To this end, we formalize Dynamic Logic as a set of locales which build upon
each other and which can be instantiated for arbitrary dynamic logics that
satisfy the constraints imposed by the locales' assumptions.

We hope that this can help the developers of dynamic logics in multiple regards:
First, by implementing their dynamic logic based on the Heterogeneous Dynamic Logic package,
they can focus on the particularities of their dynamic logic while getting many
of the well known axioms and properties of dynamic logics ``for free''.
Secondly, due to the common core of assumptions across dynamic logics,
this package allows for the construction of \emph{Heterogenous Dynamic Logics}
(also called \emph{Heterogeneous Dynamic Theories})
which allow the (often relatively complete)
combination of multiple program logics.
For a more elaborate introduction to the theory behind Heterogeneous Dynamic Logic,
the interested reader is referred to the paper~\todo{cite}.
This document will focus on the particularities of the Isabelle formalization.

The remainder of this document is structured as follows:
We begin by formalizing first-order logic (\todo{ref}) as a locale.
Subsequently, we built upon this locale to formalize a locale for a dynamic logic (\todo{ref}).
Here we prove the first dynamic logic axioms as well as the coincidence and bounded effect lemmas.
Subsequently, we introduce \emph{lifting} which allows us to take one dynamic logic locale and extend
the functionality of its programs.
We demonstrate this on two examples (havoc operator and regular closure).
We also introduce a futher constrained locale to reason about loop convergence.
Finally, we introduce \emph{Heterogeneous Dynamic Theories} as a means to combine multiple dynamic logic locales into a
single dynamic logic.

\tableofcontents

% sane default for proof documents
\parindent 0pt\parskip 0.5ex

% generated text of all theories
\input{session}

% optional bibliography
\bibliographystyle{abbrv}
\bibliography{root}

\end{document}

%%% Local Variables:
%%% mode: latex
%%% TeX-master: t
%%% End:
